% GENERATED FILE. Edit canonical lessons, then run npm run generate:books.
% canonical-source-hash: fnv1a64:eb1a14ea774982ae
% canonical-lessons: ZH-C493-buxi, ZH-C493-jiajiao, ZH-W493-jiao-instruct, ZH-C493-jiaoyu, ZH-C493-fudao, ZH-C493-qingjiao

\chapter{Extra Lessons}
\label{ch:zh-extra-lessons}
\hlchaptermodality{493}

\begin{chapteropening}
\textbf{By the end of this chapter:} \emph{I can talk about extra lessons, a private tutor, education, tutoring and asking for advice.}

Complete the last lesson of chapter 493: I can talk about extra lessons, a private tutor, education, tutoring and asking for advice.
\end{chapteropening}

\section[bǔxí]{bǔxí — to take extra lessons}

\label{lesson:ZH-C493-buxi}



\begin{quote}
Before the new one: say the Mandarin for English (the language), then the Mandarin for a major (at university).
\end{quote}

\subsection*{You'll want to know: bǔxí}
\emph{bǔxí} — \textquotedblleft{}to take extra lessons\textquotedblright{}.

\begin{quote}
\emph{wǒ bǔxí le yì nián shùxué} — I had extra maths lessons for a year.
\end{quote}

Its characters come later; for now it is a word for the ear.

\begin{cousinweb}
\emph{bǔxí} is built from \emph{bǔ} 'to patch' and \emph{xí} 'to practise', both inherited from Middle Chinese, \emph{puX} and \emph{zip} (Guangyun, in Baxter's transcription), which glosses the first 'to mend clothes' and the second 'to learn'.
\end{cousinweb}

\begin{grammarlens}[title={What you've built}]
One more word for saying how it was.
\end{grammarlens}

\subsection*{Your turn}
\begin{itemize}
\raggedright
  \item \emph{Say it:} \emph{bǔxí}
  \item \emph{Say it:} \emph{bǔxí}, once more
  \item \emph{Say it:} \emph{wǒ bǔxí le yì nián shùxué} — I had extra maths lessons for a year.
  \item \emph{Recall:} say the Mandarin for English (the language), then the Mandarin for a major (at university), then say \emph{bǔxí} again
\end{itemize}

\begin{tcolorbox}[breakable,colback=teal!4,colframe=teal!35!black,title={Before you move on}]
What does bǔxí mean? (\textquotedblleft{}To take extra lessons\textquotedblright{}.) Say it once more.
\end{tcolorbox}
\section[\texorpdfstring{\zh{家教} (jiājiào)}{家教 (jiājiào)}]{\zh{家教} (jiājiào) — private tutor}

\label{lesson:ZH-C493-jiajiao}



\begin{quote}
Before the new one: say the Mandarin for a major (at university), then the Mandarin for to take extra lessons.
\end{quote}

\subsection*{You'll want to know: \zh{家教}}
\textbf{\zh{家教}} — \emph{jiājiào} — \textquotedblleft{}private tutor\textquotedblright{}.

\begin{quote}
\textbf{\zh{大学的时候我当过家教}} \emph{dàxué de shíhou wǒ dāng guo jiājiào} — At university I worked as a private tutor.
\end{quote}

The new character, in pieces: \zh{教} — eleven strokes in two pieces: on the left the top of \zh{老} \emph{lǎo} with \zh{子} \emph{zǐ} tucked under it, seven strokes; on the right four strokes, two falling to the left, a short bar and a long stroke falling to the right. The next lesson puts it in your hand.

\begin{cousinweb}
\textbf{\zh{家教}} (\emph{jiājiào}) is built from \zh{家} 'a home' and \zh{教} 'to teach', both inherited from Middle Chinese, \emph{kae} and \emph{kaewH} (Guangyun, in Baxter's transcription), which glosses the second 'to instruct'. Here \zh{教} keeps the falling tone, jiào.
\end{cousinweb}

\begin{grammarlens}[title={What you've built}]
One more word for saying how it was.
\end{grammarlens}

\subsection*{Your turn}
\begin{itemize}
\raggedright
  \item \emph{Say it:} \emph{jiājiào}
  \item \emph{Say it:} \emph{jiājiào}, once more
  \item \emph{Say it:} \emph{dàxué de shíhou wǒ dāng guo jiājiào} — At university I worked as a private tutor.
  \item \emph{Recall:} say the Mandarin for a major (at university), then the Mandarin for to take extra lessons, then say \emph{jiājiào} again
\end{itemize}

\begin{tcolorbox}[breakable,colback=teal!4,colframe=teal!35!black,title={Before you move on}]
What does \zh{家教} mean? (\textquotedblleft{}Private tutor\textquotedblright{}.) Say it once more.
\end{tcolorbox}
\section[\texorpdfstring{\zh{教} (jiào)}{教 (jiào)}]{\zh{教} — the top of lǎo over zǐ, then four strokes}

\label{lesson:ZH-W493-jiao-instruct}



\begin{quote}
The word first, then the character inside it.

\begin{itemize}
\raggedright
  \item \emph{Recall:} say \emph{jiājiào} — \textbf{R1}, one lesson back, before it decays
\end{itemize}

\begin{quote}
\zh{老} \zh{子} \zh{家}
\end{quote}

\emph{Read it:} each one aloud, with nothing said first — \emph{lǎo}, \emph{zǐ}, \emph{jiā}
\end{quote}

\subsection*{Script — the top of lǎo over zǐ, then four strokes}
\hlblockfigure{\detokenize{figures/ZH-W493-jiao-instruct-filmstrip.pdf}}{How \zh{教} is written, stroke by stroke}

\begin{quote}
\zh{教}
\end{quote}

\textbf{Two pieces, left before right.}

\textbf{The left piece, seven strokes.} First the top of \zh{老}, four strokes, as you wrote them: a short bar, a vertical through it, a long bar, and a long stroke falling to the left. Then \zh{子} \emph{zǐ}, three strokes, tucked under it: a bar that turns and falls to the left, a vertical with a hook, and a stroke rising to the right.

\textbf{The right piece, four strokes}: a short stroke falling to the left, a bar, a long stroke falling to the left, and a long stroke falling to the right.

4 + 3 + 4 = eleven.

\subsection*{Your turn}
\begin{itemize}
\raggedright
  \item \emph{Write it:} \zh{教}, all eleven, with the model in front of you
  \item \emph{Cover:} write it again from memory, and count the pieces, not the strokes
\end{itemize}

\begin{tcolorbox}[breakable,colback=teal!4,colframe=teal!35!black,title={Before you move on}]
How many pieces does \zh{教} have? (\textbf{Two}.) What starts the left piece? (\textbf{The top of \zh{老}}.) What sits under it? (\textbf{\zh{子}}.)

Source: \href{https://raw.githubusercontent.com/chanind/hanzi-writer-data/68d10a4b21150cae5e1ebbd223eed289cf32d90c/data/\%E6\%95\%99.json}{Hanzi Writer Data \zh{教}, pinned stroke paths and medians}.
\end{tcolorbox}
\section[jiàoyù]{jiàoyù — to educate; to bring up (the young); education}

\label{lesson:ZH-C493-jiaoyu}



\begin{quote}
Before the new one: say the Mandarin for to take extra lessons, then the Mandarin for a private tutor.
\end{quote}

\subsection*{You'll want to know: jiàoyù}
\emph{jiàoyù} — \textquotedblleft{}to educate; to bring up (the young); education\textquotedblright{}.

\begin{quote}
\emph{jiàoyù háizi bù róngyì} — Bringing up children isn't easy.
\end{quote}

Its characters come later; for now it is a word for the ear.

\begin{cousinweb}
\emph{jiàoyù} is built from \zh{教} 'to instruct' and \emph{yù} 'to rear', both inherited from Middle Chinese, \emph{kaewH} and \emph{yuwk} (Guangyun, in Baxter's transcription), which glosses the first 'to instruct' and the second 'to nourish' and 'to grow'.
\end{cousinweb}

\begin{grammarlens}[title={What you've built}]
One more word for saying how it was.
\end{grammarlens}

\subsection*{Your turn}
\begin{itemize}
\raggedright
  \item \emph{Say it:} \emph{jiàoyù}
  \item \emph{Say it:} \emph{jiàoyù}, once more
  \item \emph{Say it:} \emph{jiàoyù háizi bù róngyì} — Bringing up children isn't easy.
  \item \emph{Recall:} say the Mandarin for to take extra lessons, then the Mandarin for a private tutor, then say \emph{jiàoyù} again
\end{itemize}

\begin{tcolorbox}[breakable,colback=teal!4,colframe=teal!35!black,title={Before you move on}]
What does jiàoyù mean? (\textquotedblleft{}To educate; to bring up (the young); education\textquotedblright{}.) Say it once more.
\end{tcolorbox}
\section[fǔdǎo]{fǔdǎo — to tutor; to give guidance}

\label{lesson:ZH-C493-fudao}



\begin{quote}
Before the new one: say the Mandarin for a private tutor, then the Mandarin for to educate.
\end{quote}

\subsection*{You'll want to know: fǔdǎo}
\emph{fǔdǎo} — \textquotedblleft{}to tutor; to give guidance\textquotedblright{}.

\begin{quote}
\emph{jiějie měi tiān wǎnshang fǔdǎo wǒ zuò zuòyè} — My sister helps me with my homework every evening.
\end{quote}

Its characters come later; for now it is a word for the ear.

Two third tones meet in this word: fǔ is said with a rising tone before dǎo, and the spelling keeps both marks.

\begin{cousinweb}
\emph{fǔdǎo} is built from \emph{fǔ} 'to help' and \emph{dǎo} 'to lead', both inherited from Middle Chinese, \emph{bjuX} and \emph{dawH} (Guangyun, in Baxter's transcription), which glosses the first 'to assist' and the second 'to guide'. Both modern third tones are irregular: \emph{bjuX} and \emph{dawH} would regularly give fù and dào.
\end{cousinweb}

\begin{grammarlens}[title={What you've built}]
One more word for saying how it was.
\end{grammarlens}

\subsection*{Your turn}
\begin{itemize}
\raggedright
  \item \emph{Say it:} \emph{fǔdǎo}
  \item \emph{Say it:} \emph{fǔdǎo}, once more
  \item \emph{Say it:} \emph{jiějie měi tiān wǎnshang fǔdǎo wǒ zuò zuòyè} — My sister helps me with my homework every evening.
  \item \emph{Recall:} say the Mandarin for a private tutor, then the Mandarin for to educate, then say \emph{fǔdǎo} again
\end{itemize}

\begin{tcolorbox}[breakable,colback=teal!4,colframe=teal!35!black,title={Before you move on}]
What does fǔdǎo mean? (\textquotedblleft{}To tutor; to give guidance\textquotedblright{}.) Say it once more.
\end{tcolorbox}
\section[\texorpdfstring{\zh{请教} (qǐngjiào)}{请教 (qǐngjiào)}]{\zh{请教} (qǐngjiào) — to seek advice (from someone who knows more)}

\label{lesson:ZH-C493-qingjiao}



\begin{quote}
Before the new one: say the Mandarin for to educate, then the Mandarin for to tutor.
\end{quote}

\subsection*{You'll want to know: \zh{请教}}
\textbf{\zh{请教}} — \emph{qǐngjiào} — \textquotedblleft{}to seek advice (from someone who knows more)\textquotedblright{}.

\begin{quote}
\textbf{\zh{我常常请教我的上司}} \emph{wǒ chángcháng qǐngjiào wǒ de shàngsi} — I often ask my boss for advice.
\end{quote}

\begin{cousinweb}
\textbf{\zh{请教}} (\emph{qǐngjiào}) is built from \zh{请} 'to ask' and \zh{教} 'to instruct', both inherited from Middle Chinese, \emph{tshjengX} and \emph{kaewH} (Guangyun, in Baxter's transcription), which glosses the first 'to beg' and 'to ask'.
\end{cousinweb}

\begin{grammarlens}[title={What you've built}]
One more word for saying how it was.
\end{grammarlens}

\subsection*{Your turn}
\begin{itemize}
\raggedright
  \item \emph{Say it:} \emph{qǐngjiào}
  \item \emph{Say it:} \emph{qǐngjiào}, once more
  \item \emph{Say it:} \emph{wǒ chángcháng qǐngjiào wǒ de shàngsi} — I often ask my boss for advice.
  \item \emph{Recall:} say the Mandarin for to educate, then the Mandarin for to tutor, then say \emph{qǐngjiào} again
\end{itemize}

\begin{tcolorbox}[breakable,colback=teal!4,colframe=teal!35!black,title={Before you move on}]
What does \zh{请教} mean? (\textquotedblleft{}To seek advice (from someone who knows more)\textquotedblright{}.) Say it once more.
\end{tcolorbox}
